% =========================================================
% Volume IV  \S{}1.7.?  --- Z as a Model of the Arithmetic Signature
% WIRING NOTE (cross-volume):
%   Verify exact Volume II labels before final commit.
%
% Expected targets:
%   def:arithmetic-signature
%   def:addition-on-z-classes, def:multiplication-on-z-classes
%   def:algebraic-ring
%   def:positive-cone
%   thm:z-is-ring
%   thm:z-is-integral-domain
%   thm:z-is-ordered-ring
%   thm:n-embeds-in-z
%   thm:z-embeds-in-q
%   def:structure-expansion
%   def:structure-reduction
% =========================================================

\subsection*{The Integers as a Structure}

\begin{remark*}[Modeling Goal]
We exhibit $\mathbb{Z}$ as a concrete \emph{model}: fix the ring signature,
take the carrier and operations constructed in Volume~II, interpret each
symbol, and discharge satisfaction by citing the Volume~II theorem that
$\mathbb{Z}$ satisfies the ring axioms. Order is then added by expansion, and
various algebraic fragments are recovered by reduction.
\end{remark*}

% ---------------------------------------------------------
% Signature
% ---------------------------------------------------------

\begin{remark*}[Exposition]
A model is a signature together with an interpretation on a carrier. The
signature here is the ring fragment of the
\hyperref[def:arithmetic-signature]{arithmetic signature},
\[
\mathcal{L}_{\mathrm{ring}}
=
\{+,\cdot,-,0,1\},
\]
with binary function symbols $+,\cdot$, unary function symbol $-$, and
constant symbols $0,1$.

The order symbol $<$ is deliberately absent. It enters only through expansion.
\end{remark*}

% ---------------------------------------------------------
% Model
% ---------------------------------------------------------

\begin{definition}[The Integer Model $\mathcal{Z}$]
\label{def:integer-model}

The \textbf{integer model} is the
$\mathcal{L}_{\mathrm{ring}}$-structure

\[
\mathcal{Z}
=
\bigl(
\mathbb{Z};
+^{\mathcal Z},
\cdot^{\mathcal Z},
-^{\mathcal Z},
0^{\mathcal Z},
1^{\mathcal Z}
\bigr),
\]

whose carrier $|\mathcal Z|=\mathbb Z$ is the set of integers constructed in
Volume~II, and whose symbols are interpreted by the
\hyperref[def:addition-on-z-classes]{integer addition} and
\hyperref[def:multiplication-on-z-classes]{integer multiplication}.

Thus

\[
+^{\mathcal Z}
\]

is integer addition,

\[
\cdot^{\mathcal Z}
\]

is integer multiplication,

\[
-^{\mathcal Z}
\]

is integer negation, and

\[
0^{\mathcal Z},
\quad
1^{\mathcal Z}
\]

are the distinguished integer constants.
\end{definition}

% ---------------------------------------------------------
% Quantified statement
% ---------------------------------------------------------

\begin{remark*}[Standard quantified statement]

\[
|\mathcal Z|
=
\mathbb Z,
\]

\[
+^{\mathcal Z},
\cdot^{\mathcal Z}
:
\mathbb Z\times\mathbb Z
\to
\mathbb Z,
\]

\[
-^{\mathcal Z}
:
\mathbb Z
\to
\mathbb Z,
\]

\[
0^{\mathcal Z},
1^{\mathcal Z}
\in
\mathbb Z.
\]

\end{remark*}

% ---------------------------------------------------------
% Dependencies
% ---------------------------------------------------------% ---------------------------------------------------------
% Interpretation
% ---------------------------------------------------------

\begin{remark*}[Interpretation]

Nothing new is constructed here.

The carrier and operations were already built in Volume~II.

The purpose of the model

\[
\mathcal Z
\]

is to package the integer system as a single mathematical object so that
theorems may quantify over rings and then be applied directly to the integers.

The signature is the type.

The structure

\[
\mathcal Z
\]

is the instance.

\end{remark*}

% ---------------------------------------------------------
% Satisfaction
% ---------------------------------------------------------

\begin{remark*}[Satisfaction]

The structure

\[
\mathcal Z
\]

is a model of the ring axioms:

\[
\mathcal Z
\models
\Sigma_{\mathrm{ring}},
\]

equivalently,

\[
\operatorname{Ring}
(\mathbb Z,+,\cdot,0,1).
\]

This is not re-proved here.

It is exactly the content of

\hyperref[thm:z-is-ring]{$\mathbb Z$ Is a Ring}

in Volume~II, which verifies each ring axiom against the integer operations.

Likewise,

\[
\mathcal Z
\models
\Sigma_{\mathrm{id}},
\]

because

\hyperref[thm:z-is-integral-domain]{$\mathbb Z$ Is an Integral Domain}

shows that multiplication on the integers has no zero divisors.

\end{remark*}

% ---------------------------------------------------------
% Expansion
% ---------------------------------------------------------

\begin{remark*}[Expansion to the Ordered Ring]

Adjoining the relation symbol $<$ gives the expansion

\[
\mathcal Z_{<}
=
(\mathbb Z;
+,
\cdot,
-,
0,
1,
<),
\]

\[
\operatorname{Expansion}
(\mathcal Z_{<},\mathcal Z),
\]

in the sense of
\hyperref[def:structure-expansion]{Structure Expansion}.

The carrier is unchanged and exactly one interpreted symbol is added.

The relation

\[
<^{\mathcal Z}
\]

is the usual integer order determined by the
\hyperref[def:positive-cone]{positive cone}.

That this expansion satisfies the ordered-ring axioms is witnessed by

\hyperref[thm:z-is-ordered-ring]{$\mathbb Z$ Is an Ordered Ring}

in Volume~II.

\end{remark*}

% ---------------------------------------------------------
% Additive reduct
% ---------------------------------------------------------

\begin{remark*}[Reduct to the Additive Group]

Forgetting

\[
\cdot,
1
\]

leaves the reduct

\[
\mathcal Z
\!\restriction\!
\{+,-,0\}
=
(\mathbb Z;+,-,0),
\]

\[
\operatorname{Reduction}
\!\bigl(
\mathcal Z
\!\restriction\!
\{+,-,0\},
\mathcal Z
\bigr).
\]

This reduct is the additive abelian group of integers.

It remembers addition and additive inverses while ignoring multiplication.

\end{remark*}

% ---------------------------------------------------------
% Multiplicative reduct
% ---------------------------------------------------------

\begin{remark*}[Reduct to the Multiplicative Monoid]

Forgetting

\[
+,
-,
0
\]

leaves

\[
\mathcal Z
\!\restriction\!
\{\cdot,1\}
=
(\mathbb Z;\cdot,1),
\]

the multiplicative commutative monoid of integers.

\end{remark*}

% ---------------------------------------------------------
% Blueprint
% ---------------------------------------------------------

\begin{remark*}[Blueprint Position]

The integer model occupies the next stage after the natural numbers in the
number-system hierarchy.

It sits on the embedding chain

\[
\mathbb N
\hookrightarrow
\mathbb Z
\hookrightarrow
\mathbb Q
\hookrightarrow
\mathbb R.
\]

The first embedding is the canonical map

\[
n \longmapsto n,
\]

witnessed by

\hyperref[thm:n-embeds-in-z]{$\mathbb N$ Embeds in $\mathbb Z$}.

The second embedding is witnessed by

\hyperref[thm:z-embeds-in-q]{$\mathbb Z$ Embeds in $\mathbb Q$}.

\end{remark*}

% ---------------------------------------------------------
% Structural meaning
% ---------------------------------------------------------

\begin{remark*}[Interpretation]

The passage

\[
\mathbb N
\longrightarrow
\mathbb Z
\]

supplies exactly one new algebraic capability:

\[
\forall a\in\mathbb Z,
\quad
\exists(-a)\in\mathbb Z.
\]

Every integer has an additive inverse.

This is the structural feature absent from the natural numbers.

Consequently,

\[
(\mathbb N,+)
\]

is only a commutative monoid, whereas

\[
(\mathbb Z,+)
\]

is an abelian group.

The passage

\[
\mathbb Z
\longrightarrow
\mathbb Q
\]

will later supply multiplicative inverses for nonzero elements and thereby
upgrade the integer ring into a field.

\end{remark*}

% ---------------------------------------------------------
% Closing observation
% ---------------------------------------------------------

\begin{remark*}[Structural Significance]

The integer model is the first standard number system that satisfies the full
ring signature.

The natural numbers support addition and multiplication, but not additive
inverses.

The integers close precisely that gap.

For this reason $\mathbb Z$ occupies the ring layer of the structural hierarchy,
sitting between the semiring structure of $\mathbb N$ and the field structure of
$\mathbb Q$.

\end{remark*}

\begin{dependencies}

\begin{itemize}
    \item \hyperref[def:arithmetic-signature]{Arithmetic Signature}.
    \item \hyperref[def:algebraic-ring]{Ring}.
    \item \hyperref[def:integers]{The Integers}.
    \item \hyperref[def:addition-on-z-classes]{Integer Addition}.
    \item \hyperref[def:multiplication-on-z-classes]{Integer Multiplication}.
\end{itemize}

\end{dependencies}
